\documentclass[12pt]{article}
\input{iati_structure.tex}
\renewcommand{\bottomtitlespace}{3cm}

\usepackage[english]{babel}

\begin{document}

\renewcommand{\headerLang}{English}
\renewcommand{\problemName}{AB13. Vision}

\renewcommand{\headerLeft}{IATI Day 1, Senior group}
\renewcommand{\headerRightFirst}{XVII INTERNATIONAL ADVANCED TOURNAMENT IN INFORMATICS}
\renewcommand{\headerRightSecond}{ONLINE 2026}

\renewcommand{\tl}{$5$ s}
\renewcommand{\ml}{$1024$ MB}
\problem{Задача \problemName}
\addauthor{Автор: Emil Indzhev}{2026}{04}{16}

Вселенскиот брод USS Enterprise патува низ универзумот, кој претставува бесконечна 2D мрежа од сектори. Во секој сектор $(X, Y)$\footnote{$X$ е број на редица (се зголемува надолу), а $Y$ е број на колона (се зголемува надесно).} има светилник со јачина на видливост $V_{X, Y}$. Кога се наоѓа во тој сектор, бродот го користи светилникот за да ги скенира сите сектори $(X', Y')$ за кои важи $X - V_{X, Y} \leq X' \leq X + V_{X, Y}$ и  $Y - V_{X, Y} \leq Y' \leq Y + V_{X, Y}$. Така, видените сектори формираат квадрат со должина на страната $2V_{X, Y} + 1$. За секој таков сектор, единствената информација што се гледа е јачината на видливост на светилникот во него, односно $V_{X', Y'}$ (ова е затоа што светилниците дејствуваат и како предаватели и како приемници). По скенирањето, бродот потоа се телепортира во еден од овие видени сектори, бидејќи светилниците дејствуваат и како телепортери.

За жал, светилниците имаат една критична маана: тие може да ја превртат $X$-оската на скенирањето, или $Y$-оската, или двете (односно постојат $4$ можни ориентации на скенирањето). Забележете дека бродот го користи ова евентуално превртено скенирање за да одлучи во кој сектор да се телепортира, но и телепортацијата ја користи истата ориентација, на пример ако $Y$-оската е превртена и бродот се телепортира во левата насока (намалување на $Y$) на скенирањето, тој всушност ќе се телепортира надесно. Ова значи дека бродот навистина ќе се телепортира во секторот избран на скенирањето.

Понатаму, бродот нема меморија, па неговата одлука за тоа каде да се телепортира мора да зависи само од скенирањето произведено од светилникот. Дополнително, стратегијата на бродот мора да биде детерминистичка -- ако добие исто скенирање, мора да избере да се телепортира во истиот сектор од скенирањето.

Бродот започнува на непознати координати и мора секогаш (без разлика од каде започнува и без разлика какви ориентации на скенирањата добива) да може да патува во насока на зголемување на $X$ и $Y$ произволно далеку (можеме да кажеме дека бараме на крај да достигне сектор $(X, Y)$ таков што $X, Y \geq 10^{100}$). Сепак, за ова да биде можно, клучни се и стратегијата на патување на бродот и распоредот на јачините на видливост низ универзумот.

Тука настапуваш ти -- треба да дизајнираш правоаголна шема $P$ со димензии $N$ на $M$ од јачини на видливост, која бесконечно ќе се повторува низ универзумот: $V_{X, Y} = P_{X \bmod N, Y \bmod M}$. Исто така треба да ја дизајнираш стратегијата на патување на бродот, односно функција што го прима скенирањето произведено од светилникот и избира во кој од тие сектори бродот да се телепортира понатаму.

Бидејќи светилниците со поголема јачина на видливост се скапи, твојата цел е да произведеш валидно решение, притоа обидувајќи се да ја минимизираш просечната јачина на видливост во твојата шема, односно просечната вредност на $P$.

Бидејќи овој проблем може да изгледа доста предизвикувачки, решаваш да започнеш со задача за вежбање: истиот проблем, но во 1D. Во оваа поедноставна задача, ја отфрламе $X$ димензијата и ја користиме само $Y$. Така, твојата шема е само низа со должина $M$ и таа се повторува вака: $V_Y = P_{Y \bmod M}$. Скенирањата произведени од светилниците исто така се само низи со должина $2V_Y + 1$ (кои ги содржат секторите $Y'$ такви што  $Y - V_{X, Y} \leq Y' \leq Y + V_{X, Y}$). Тука постојат само 2 можни ориентации за скенирањето (обично или превртено), а бродот мора да може да достигне $Y \geq 10^{100}$.

Твоите решенија за 1D случајот и за 2D случајот ќе се бодуваат одделно и може да добиеш поени за едниот без да имаш валидно решение за другиот.

\subsection{Детали за имплементација}

Треба да имплементираш 4 функции: \verb|getVisionPattern1d|, \verb|getMove1d|, \verb|getVisionPattern2d| и \verb|getMove2d|. Првите 2 се користат за 1D случајот, а вторите 2 -- за 2D случајот. Во секој тест ќе биде повикан само еден од овие парови функции, но сите 4 мора да бидат имплементирани (дури и ако имплементациите за едниот пар се празни) за решението да биде валидно (да се компајлира).

\begin{lstlisting}
 std::vector<int> getVisionPattern1d()
\end{lstlisting}

Оваа функција ќе биде повикана еднаш, ако тестот е за 1D случајот. Таа мора да ја врати низата $P$ со должина $M$ -- 1D шемата на јачини на видливост што ќе се повторува. $M$ и сите елементи на $P$ мора да бидат меѓу $1$ и $60$.

\begin{lstlisting}
 int getMove1d(std::vector<int> v)
\end{lstlisting}

Оваа функција ќе биде повикана еднаш за секое можно скенирање, ако тестот е за 1D случајот. Таа добива скенирање произведено од светилник, како низа со должина $2V + 1$ (каде што $V$ е јачината на видливост на светилникот во централниот сектор од низата). Таа мора да врати во кој сектор бродот треба да се телепортира, како индекс $T$ во низата, така што $0 \leq T < 2V + 1$.

\begin{lstlisting}
 std::vector<std::vector<int>> getVisionPattern2d()
\end{lstlisting}

Оваа функција ќе биде повикана еднаш, ако тестот е за 2D случајот. Таа мора да врати правоаголна табела $P$ со димензии $N$ на $M$ -- 2D шемата на јачини на видливост што ќе се повторува. $N$, $M$ и сите елементи на $P$ мора да бидат меѓу $1$ и $60$.

\begin{lstlisting}
 std::pair<int, int> getMove2d(std::vector<std::vector<int>> v)
\end{lstlisting}

Оваа функција ќе биде повикана еднаш за секое можно скенирање, ако тестот е за 2D случајот. Таа добива скенирање произведено од светилник, како квадратна табела со димензии $2V + 1$ на $2V + 1$ (каде што $V$ е јачината на видливост на светилникот во централниот сектор од квадратот). Таа мора да врати во кој сектор бродот треба да се телепортира, како пар индекси $S$ и $T$ во табелата, така што $0 \leq S, T < 2V + 1$.

За дадената димензионалност $D$ ($1$ или $2$) на тестот, оценувачот ќе провери дали твојата програма произведува валидна шема -- дали прави валидни избори за телепортација за сите можни скенирања и дали твоето решение работи без разлика на почетните координати и низата ориентации на скенирањата.

\subsection{Ограничувања}

\vspace{0.1em}

\begin{itemize}
\item $1 \leq D \leq 2$
\item $1 \leq N, M, V_{X, Y}, V_Y \leq 60$
\end{itemize}

\newpage

\subsection{Подзадачи и бодување}

\begin{table}[H]
\begin{tblr}{|Q[c,m]|Q[c,m]|Q[c,m]|Q[c,m]|}
\hline
\textbf{Подзадача} & \textbf{Поени} & \textbf{$D$} & \textbf{$A^*$} \\
\hline
$1$ & $24$ & $1$ & $1.5$ \\
\hline
$2$ & $76$ & $2$ & $1.125$ \\
\hline
\end{tblr}
\end{table}
\FloatBarrier

Твојот резултат за дадена подзадача (ако решението за неа ти е точно) зависи од просечната јачина на видливост во твојата шема $P$ за неа. Нека тоа е $A$, а нека целната вредност за подзадачата е $A^*$. Твојот резултат (дел од поените за подзадачата) ќе биде:

$$1 - {\left(1 - \frac{A^* - 1}{\max{\left(A, A^*\right)} - 1}\right)}^{0.8}$$

\subsection{Пример шема}

Да претпоставиме дека твоето решение ја враќа следната шема со $N=3$ и $M=2$:

\[
\left[
\begin{array}{c c}
1 & 2 \\
3 & 4 \\
5 & 6
\end{array}
\right]
\]

Сега, да ги разгледаме можните скенирања во секторот $(0, 1)$, кој има јачина на видливост $2$. Постојат $4$ можни ориентации (од кои некои го произведуваат истото скенирање на оваа шема):

\[
\begin{array}{cc}
\text{Orientation $0$: No flips} & \text{Orientation $1$: $Y$-axis flipped} \\
\left[
\begin{array}{ccccc}
4 & 3 & 4 & 3 & 4 \\
6 & 5 & 6 & 5 & 6 \\
2 & 1 & 2 & 1 & 2 \\
4 & 3 & 4 & 3 & 4 \\
6 & 5 & 6 & 5 & 6
\end{array}
\right]
&
\left[
\begin{array}{ccccc}
4 & 3 & 4 & 3 & 4 \\
6 & 5 & 6 & 5 & 6 \\
2 & 1 & 2 & 1 & 2 \\
4 & 3 & 4 & 3 & 4 \\
6 & 5 & 6 & 5 & 6
\end{array}
\right]
\\
\\
\text{Orientation $2$: $X$-axis flipped} & \text{Orientation $3$: Both axes flipped} \\
\left[
\begin{array}{ccccc}
6 & 5 & 6 & 5 & 6 \\
4 & 3 & 4 & 3 & 4 \\
2 & 1 & 2 & 1 & 2 \\
6 & 5 & 6 & 5 & 6 \\
4 & 3 & 4 & 3 & 4
\end{array}
\right]
&
\left[
\begin{array}{ccccc}
6 & 5 & 6 & 5 & 6 \\
4 & 3 & 4 & 3 & 4 \\
2 & 1 & 2 & 1 & 2 \\
6 & 5 & 6 & 5 & 6 \\
4 & 3 & 4 & 3 & 4
\end{array}
\right]
\end{array}
\] 

Бидејќи $0$ и $1$ се идентични и бидејќи бродот е детерминистички, за нив ќе биде направен само еден повик до \verb|getMove2d|. Да претпоставиме дека твоето решение го враќа потегот $(0, 1)$ за ова скенирање. Кај ориентацијата $0$ ова би значело движење $1$ налево и $2$ нагоре, додека кај ориентацијата $1$ -- движење $1$ надесно и $2$ нагоре.

\subsection{Пример оценувач}

Првиот влез на пример-оценувачот е $D$. После тоа тој ќе ги направи сите повици до твојата функција: прво ќе ја добие шемата $P$, а потоа ќе ги кешира сите избори за телепортација. После тоа, ќе започне интеракција преку конзола, барајќи ги почетните координати на бродот. Потоа, ќе ја испечати тековната состојба: координатите и ``суровото'' (без превртувања) скенирање. Потоа, ќе ја побара ориентацијата што треба да се користи ($0$ за без превртувања, $1$ за превртување на $Y$-оската, $2$ за превртување на $X$-оската и $3$ за превртување на двете). После тоа, ќе го испечати скенирањето во дадената ориентација, како и кој потег го избрал бродот. Бродот ќе биде поместен и ова ќе се повторува. Забележи дека пример-оценувачот нема автоматски да провери дали твоето решение е точно.

\end{document}